\section{Einleitung}

Das Konzept der gekoppelten Pendel ist ein wichtiges Konzept in der Molekülphysik und Molekülchemie, da dadurch die Bewegungen von Molekülen angenähert werden können. Der Versuch ist der Mechanik unterzuordnen, damit sind die verwendeten Messmethoden zum einen Meterstäbe und Skalen, um die Auslenkung der Pendel zu ermitteln. Die andere Messmethode ist eine Stoppuhr, um Frequenzen und Schwingungsperioden zu bestimmen.
